\documentclass[11pt]{book}
\usepackage[utf8]{inputenc}
\usepackage{geometry}
\geometry{a5paper, margin=1.5cm}
\usepackage{ebgaramond}
\usepackage{epigraph}
\usepackage{verse}
\usepackage{microtype}
\setlength{\epigraphwidth}{0.7\textwidth}
\setlength{\epigraphrule}{0pt}

\begin{document}

\chapter*{The Justicar's Judgment}
\addcontentsline{toc}{chapter}{The Justicar's Judgment}

\vspace{2cm}
\begin{center}
\textit{A tale of Viterra, told by a clerk who sharpened the quills}
\end{center}
\vspace{1cm}

\epigraph{The law does not forget. Neither does the son who was not ransomed.}{--- Queen's Justicar Elara Venn, from the bench}

\section*{The Trial of Baron Aldric the Turncloak}

\textit{Set down by Dusana of the Raven Road, after a cup of small beer in a Viterran hedge-court. The ink is gall and the memory of a father's silence.}

\subsection*{The Charges}

In the fifth year of Queen Veralyn's peace, the Queen's Justicar---a woman named Elara Venn---received a sealed writ. The wax was black, stamped with the crown of Viterra and the broken column of Vhasia. Inside, a single sheet of vellum:

\emph{You are commanded to hear the case of Baron Aldric of the Scar, formerly Sir Aldric of the White Moor, accused of treason thrice sworn and thrice forsworn, of oath-breaking before witnesses, and of refusing the lawful ransom of his own blood. Let the ledger be balanced.}

Elara read the name. Her hand did not shake.

She had squired for Aldric twenty years ago, when he was a young knight with a quick laugh and a quicker sword. She had carried his shield at the Ford of Ashes. She had watched him bargain for his first holding---a burnt tower and a handful of serfs---with the same smile he used to haggle over a horse.

She had also watched him abandon Queen Veralyn's mother at the Battle of the Three Bridges, switching his banner to the Pretender's side an hour before the charge. When the Pretender lost, Aldric swore fealty again. Elara had testified on his behalf at the first inquiry.

She was young then. She believed in second chances.

\subsection*{The Knight's Ledger}

The court convened in the White Hall of Viterra, a long room lined with oak and the painted faces of dead magistrates. Baron Aldric was brought in chains of gilded bronze---a courtesy, because he still had friends. His hair had gone grey. His eyes had not.

Elara read the first count: \textit{That you swore loyalty to Queen Veralyn's father, then to the usurper Caster, then again to the Queen's father, then to the Duke of Adria, and finally to the Queen herself, each time breaking the previous oath without payment of blood-silver.}

Aldric smiled. ``A man must survive, Justicar. The throne changed hands. I changed my hand. Is that not the custom of Viterra?''

The hall murmured. Elara did not smile.

``The custom of Viterra,'' she said, ``is that an oath is a bridge. You do not burn a bridge behind you unless you have paid the toll in full. You owe eighteen years of back-tolls, Aldric. That is not survival. That is usury.''

She let the silence stretch. Then she read the second count---the one that had kept her awake.

\subsection*{The Son Who Was Not Ransomed}

``In the year of the Broken Crown, after the Battle of the Red Ditch, your son---Ser Roran, then seventeen---was captured by the forces of the Duke of Adria. The customary ransom for a knight's eldest son was three hundred silver marks. The Duke's herald offered you a discount of fifty marks if you paid within the week.''

Aldric's smile thinned. ``I recall.''

``You refused. You said, and I quote the herald's sworn testimony: 'I have more sons and can make further. The boy was weak. Let the Duke keep him.'''

A rustle passed through the hall. Elara did not look at the gallery, where a lean, grey-faced man sat in the shadows---Ser Roran, now thirty-seven, his left hand sheathed in a gauntlet of silvered steel, his right hand resting on the pommel of a sword that had seen twenty campaigns.

``The Duke did not keep him,'' Elara continued. ``He threw the boy into a debtor's cell, expecting you to relent. You did not. For three years, Roran rotted. Then, a stranger paid the ransom---a hedge-knight who had once served under your banner and remembered the boy's courage. The knight gave every coin he had. He died of a fever the following winter. Roran took his name and his sword.''

She paused.

``He did not return to you. He fought in the eastern wars. He won the Queen's own pennon at the Siege of the Horn. He knighted himself with a broken blade and the blessing of a battlefield priest. He has not spoken your name in seventeen years.''

Aldric's jaw tightened. ``The boy survived. Stronger for it.''

Ser Roran stood. The gallery fell silent. He walked to the bar and placed his gauntleted hand on the rail.

``Father,'' he said, ``do you know what I remember most? Not the cell. Not the chains. The moment I heard you say it. I was hiding behind the herald's tent. I heard you laugh.''

\subsection*{The Justicar's Judgment}

The law of Viterra is ancient and precise. Treason is punishable by death. Oath-breaking forfeits all lands. Refusing a son's ransom is not a crime---it is merely cruelty.

Elara Venn sat on the bench for three hours. She read the old charters. She consulted the \textit{Codex of Precedents}, a book so heavy that two clerks carried it. She asked for a cup of water. She asked for silence.

When she spoke, her voice was level.

``Baron Aldric of the Scar, on the first count---treason by repeated oath-breaking---I find you guilty. The sentence is forfeiture of all lands, titles, and pensions. You will be stripped of your banner and your ring. You will live as a hedge-knight if any will take you.''

Aldric's face did not change. He had expected this.

``On the second count,'' Elara continued, ``refusal of ransom and abandonment of kin---the law is silent. Viterra does not compel a father to love. But the law does have a clause for debts of the heart. It is called the \textit{Right of Acknowledgment}. No one has invoked it in two hundred years.''

She looked at Ser Roran.

``Ser Roran, son of Aldric, you have the right to demand that your father speak one truth aloud, in this court, under witness. If he refuses, he is declared \textit{anima non grata}---his soul outlawed from all Viterran burial grounds. His body will be buried at a crossroads with a stake through his chest. The law does not forbid it. The Church does not approve. But the law is the law.''

Aldric laughed. ``A truth? I have told a thousand truths. I have told the truth that a man must cut what he cannot carry. I have told the truth that sons are arrows---you loose them and you do not weep. That is my truth.''

``No,'' Ser Roran said. His voice was soft. ``Your truth is that you were afraid. Not of the Duke. Of me. I was stronger than you at seventeen. You saw it. You could not bear to pay my ransom because then you would owe me, and you have never owed anyone in your life.''

Aldric opened his mouth. Closed it.

The hall waited.

``I am not asking you to love me,'' Roran said. ``I am asking you to say it. Just once. Say: 'I was afraid.'''

Aldric stared at his son. For a long moment, the mask of the turncloak---the smile, the shrug, the armor of indifference---cracked.

``I was afraid,'' he whispered.

The sound was so small that the clerks leaned forward.

``Of what?'' Elara asked.

``Of him. Of what he would become. He would have been better than me. And I... I could not stand to watch.''

Ser Roran turned away. He did not weep. He had no tears left. He walked to the door, his silvered gauntlet clicking against the stone.

Elara raised her gavel.

``The Right of Acknowledgment is satisfied. The court is adjourned. Baron Aldric---former Baron Aldric---you are free to go. But I advise you to leave Viterra by sunset. The road is long, and the winter is coming.''

\subsection*{The Aftermath}

Aldric died three years later in a ditch outside Thepyrgos. No one claimed the body. The crossroads stake was not needed; he had outlived his own burial rights.

Ser Roran became a knight of the Queen's household. He wore the silver gauntlet over his left hand---not to hide the damage, but to remind himself that a stranger's coin had bought his life. He never took ransom. Once, a captured lord offered him five hundred marks for his freedom. Roran refused.

``Pay it to the hedge-knight's widow,'' he said. ``I have no need for gold. I had a father once. That was expensive enough.''

Queen Veralyn, when she heard the verdict, said nothing for a long time. Then she called for the \textit{Codex of Precedents} and wrote a new clause in the margin, in her own hand:

\emph{``The Right of Acknowledgment may be invoked by any child who was not ransomed, not visited, not named in a will. The truth, once spoken, is payment in full. No further interest shall accrue.''}

The clerks copied the clause into every new edition. It is still law.

\vspace{1cm}
\begin{center}
\textit{The ink is dry. The gavel is still.}
\end{center}

\end{document}
